\documentclass[12pt]{article}
\usepackage[T1]{fontenc}
\usepackage[utf8]{inputenc}
\usepackage{lmodern}
\usepackage{textcomp}
\usepackage{enumitem}
\usepackage{geometry}
\geometry{margin=1in}
\begin{document}
\begin{center}
Answer Key - Physics - Undergraduate Level Assessment 24 \
\small (Answers are ordered exactly as in the question paper.)
\end{center}

\section*{Multiple Choice}
\begin{enumerate}[label=\arabic*.]
\item \textbf{Answer:} D
\\ \textit{Options:} A) square; B) hexagon; C) cube; D) tetrahedron
\\ \textit{Note:} Correct option: D - tetrahedron
\item \textbf{Answer:} D
\\ \textit{Options:} A) Bosons have symmetric wave functions and obey the Pauli exclusion principle.; B) Bosons have antisymmetric wave functions and do not obey the Pauli exclusion principle.; C) Fermions have symmetric wave functions and obey the Pauli exclusion principle.; D) Fermions have antisymmetric wave functions and obey the Pauli exclusion principle.
\\ \textit{Note:} Correct option: D - Fermions have antisymmetric wave functions and obey the Pauli exclusion principle.
\item \textbf{Answer:} C
\\ \textit{Options:} A) There are no changes in the internal energy of the system.; B) The temperature of the system remains constant during the process.; C) The entropy of the system and its environment remains unchanged.; D) The entropy of the system and its environment must increase.
\\ \textit{Note:} Correct option: C - The entropy of the system and its environment remains unchanged.
\item \textbf{Answer:} A
\\ \textit{Options:} A) rate of change of the electric flux through S; B) electric flux through S; C) time integral of the magnetic flux through S; D) rate of change of the magnetic flux through S
\\ \textit{Note:} Correct option: A - rate of change of the electric flux through S
\item \textbf{Answer:} D
\\ \textit{Options:} A) magnetic flux through S; B) rate of change of the magnetic flux through S; C) time integral of the magnetic flux through S; D) rate of change of the electric flux through S
\\ \textit{Note:} Correct option: D - rate of change of the electric flux through S
\end{enumerate}

\section*{True / False}
\begin{enumerate}[label=\arabic*.]
\setcounter{enumi}{5}
\item \textbf{Answer:} True
\\ \textit{Options:} A) True; B) False
\\ \textit{Note:} LLM-generated true statement based on MCQ.
\item \textbf{Answer:} False
\\ \textit{Options:} A) True; B) False
\\ \textit{Note:} LLM-generated false statement. Correct answer: '4'.
\item \textbf{Answer:} True
\\ \textit{Options:} A) True; B) False
\\ \textit{Note:} LLM-generated true statement based on MCQ.
\item \textbf{Answer:} True
\\ \textit{Options:} A) True; B) False
\\ \textit{Note:} LLM-generated true statement based on MCQ.
\item \textbf{Answer:} False
\\ \textit{Options:} A) True; B) False
\\ \textit{Note:} LLM-generated false statement. Correct answer: 'interactions with the ionic lattice'.
\end{enumerate}

\section*{Long Form Response}
\begin{enumerate}[label=\arabic*.]
\setcounter{enumi}{10}
\item \textbf{Answer:} 10 eV. Explain why this is the correct answer and provide supporting evidence. Reference key facts or concepts that support this choice.
\\ \textit{Note:} Long-form derived from MCQ; award credit for stating the correct idea and giving supporting reasoning.
\item \textbf{Answer:} 2 V\_0/3. Explain why this is the correct answer and provide supporting evidence. Reference key facts or concepts that support this choice.
\\ \textit{Note:} Long-form derived from MCQ; award credit for stating the correct idea and giving supporting reasoning.
\end{enumerate}

\vfill
\noindent\textit{End of Answer Key}
\end{document}